\section{The quartic moment cone of a minimal-rank support}\label{sec:ppt3-app-moment}

We first rule out cut-product vectors in the kernel. The remaining three kernel constraints then force a unique
product ray for each qubit direction, producing the quartic curves of Theorem~\ref{thm:ppt3-universal-moments}.

\begin{lemma}[product kernel]\label{lem:ppt3-product-kernel}
\cls{proved}
A PPT operator on $\mathbb C^2\otimes\mathbb C^4$ whose kernel contains a nonzero product vector is separable. Hence the
kernel of a fully-PPT three-qubit state that is not separable across a cut contains no product vector across that cut,
and the range of every rank-five fully-PPT GME state satisfies the hypothesis of Theorem~\ref{thm:ppt3-universal-moments}.
\end{lemma}

\begin{proof}
This is a standard reduction, cf.~\cite{KCKL}; we give a short proof. Separability and PPT are invariant under local unitaries, so we may take the kernel vector to be $\ket0\ket3$. Write
$\rho=\begin{pmatrix}A&B\\B\dg&C\end{pmatrix}$ in qubit blocks. Then $\rho\ket{0,3}=0$ gives $A\ket3=0$ and $B\dg\ket3=0$.
Since $\langle0,3|\rho^{T_A}|0,3\rangle=\langle0,3|\rho|0,3\rangle=0$ and $\rho^{T_A}=\begin{pmatrix}A&B\dg\\B&C
\end{pmatrix}\succeq0$, also $\rho^{T_A}\ket{0,3}=0$, so $B\ket3=0$. Hence $\rho\ket{1,3}=\rho^{T_A}\ket{1,3}=
\ket1\otimes C\ket3$. If $C_{33}=\langle3|C|3\rangle=0$ then $C\ket3=0$ and $\rho$ lives on $\mathbb C^2\otimes\mathbb C^3$.
Otherwise put $u=C\ket3/\sqrt{C_{33}}$ and $\rho'=\rho-\ket1\bra1\otimes\ket u\bra u$. For any $M\succeq0$ and vector $f$
with $\langle f|M|f\rangle>0$, the Schur complement $M-M\ket f\bra fM/\langle f|M|f\rangle$ is positive semidefinite. With
$f=\ket{1,3}$ this applies to $M=\rho$ and to $M=\rho^{T_A}$, and both subtracted terms equal $\ket1\bra1\otimes\ket u\bra u$,
which is invariant under partial transposition. So $\rho'$ is PPT, and $\rho'\ket{j,3}=0$ for $j=0,1$: $\rho'$ lives on
$\mathbb C^2\otimes\mathbb C^3$. PPT operators on $\mathbb C^2\otimes\mathbb C^3$ are separable~\cite{HHH96}, so $\rho=\rho'+
\ket1\bra1\otimes\ket u\bra u$ is separable. For the consequence, a rank-five fully-PPT state with a product vector in
its kernel at some cut is separable across that cut, hence biseparable.
\end{proof}

\begin{proof}[Proof of Theorem~\ref{thm:ppt3-universal-moments}]
Fix a cut; after relabelling the qubits it is $A|BC$. Let $N$ be an $8\times3$ basis matrix of $R^\perp$. A vector
$a\otimes b$ lies in $R$ if and only if $M(a)b=0$, where $M(a)=N\dg(a\otimes I_4)$ is a $3\times4$ matrix depending
linearly on $a$. If $\rank M(a)<3$ for some $a\ne0$, a nonzero $y$ has $y\dg M(a)=0$, so $k=Ny\in R^\perp$ is orthogonal
to $a\otimes\mathbb C^4$, that is $k=a^\perp\otimes b'$ is a product vector, contrary to the hypothesis. So the pencil
$M(a)=a_0M_0+a_1M_1$ has rank three for every $a\ne0$.

By the Kronecker canonical form of matrix pencils~\cite[Ch.~XII]{Gant}, there are invertible $P$ and $Q$ with
$PM(a)Q$ block diagonal, with blocks $L_\varepsilon$ ($\varepsilon\times(\varepsilon+1)$, with entries $a_0$ on the
diagonal and $a_1$ on the superdiagonal), transposes $L_\eta^T$, and a regular part. Full row rank excludes every
$L_\eta^T$ block. A regular block of positive size has a homogeneous determinant that vanishes somewhere on
$\mathbb{CP}^1$, contradicting constant rank. Each remaining $L$ block contributes one kernel dimension;
nullity one therefore forces the single block $L_3$. Its kernel is spanned by
$(a_1^3,-a_0a_1^2,a_0^2a_1,-a_0^3)^T$, so $\ker M(a)$ is spanned by $b(a)=Q\,\kappa(a)$ with $\kappa(a)$ this cubic vector.
The product rays in $R$ are therefore the rays of $a\otimes b(a)=(I_2\otimes Q)(a\otimes\kappa(a))$. The eight entries of
$a\otimes\kappa(a)$ are $\pm$ quartic monomials and every one of the five monomials occurs, so $a\otimes b(a)=L\nu(a)$ for
an $8\times5$ matrix $L$ of rank five. Its column space is spanned by the curve, lies in $R$ and has dimension five, so it
equals $R$, and $L=BT_A$ for an invertible $T_A$. (This is the case of a maximal completely entangled kernel
in~\cite[Prop.~24]{CD}.) The other cuts are the same with the qubits relabelled.

A positive decomposition of an operator supported on $R$ uses only vectors in $R$, since each positive summand has zero
expectation on $R^\perp$. A biseparable $BXB\dg$ is therefore a positive combination of cut-$c$ product projectors in
$R$, that is of $BT_c\nu(a)\nu(a)\dg T_c\dg B\dg$, which gives~\eqref{eq:ppt3-universal-bs} (the converse direction is
immediate). The trace-one product projectors on $R$ form a compact set, so their cone is closed, and Carath\'eodory's
theorem in the $25$-dimensional real space of Hermitian $5\times5$ matrices bounds the number of terms. The dual
statement~\eqref{eq:ppt3-universal-dual} is the definition of the dual of a sum of cones generated by the rank-one
matrices $T_c\nu(a)\nu(a)\dg T_c\dg$.
\end{proof}
